% ProtMate 1.2.0 User Guide
%
% Compile with pdfLaTeX, XeLaTeX, LuaLaTeX or Tectonic (run twice for the
% table of contents and cross-references), for example:
%     pdflatex ProtMate_User_Guide.tex && pdflatex ProtMate_User_Guide.tex
%     tectonic ProtMate_User_Guide.tex
% Screenshots are read from the images/ folder next to this file.
%
% Copyright (C) 2024-2026 Sanjeeva Wijeyesakere. ProtMate is free software,
% released under the GNU General Public License version 3.

\documentclass[11pt,a4paper,oneside]{report}

\usepackage{iftex}
\ifPDFTeX
  \usepackage[T1]{fontenc}
  \usepackage[utf8]{inputenc}
  \usepackage{lmodern}
\else
  \usepackage{fontspec}
\fi
\usepackage[margin=2.5cm,headheight=14pt]{geometry}
\usepackage{graphicx}
\usepackage{xcolor}
\usepackage{amsmath}
\usepackage{booktabs}
\usepackage{tabularx}
\usepackage{longtable}
\usepackage{array}
\usepackage{enumitem}
\usepackage{fancyvrb}
\usepackage{fancyhdr}
\usepackage{caption}
\usepackage[most]{tcolorbox}
\usepackage[hidelinks]{hyperref}

% ---------------------------------------------------------------------------
% Settings to edit

\newcommand{\pmversion}{1.2.0}
\newcommand{\copyrightyears}{2024--2026}
\newcommand{\repourl}{https://github.com/swijeyes/ProtMate}
\newcommand{\repo}{\url{\repourl}}
\newcommand{\releases}{\url{\repourl/releases}}
\newcommand{\issues}{\url{\repourl/issues}}
% How the licence is granted; the FSF's suggested wording adds
% ", or (at your option) any later version" after "version 3 of the License".
\newcommand{\gplgrant}{either version 3 of the License}

% ---------------------------------------------------------------------------
% Appearance

\definecolor{pmblue}{HTML}{1F5FA8}
\definecolor{pmgrey}{HTML}{F3F4F6}
\definecolor{pmwarn}{HTML}{B35900}

\hypersetup{colorlinks=true,linkcolor=pmblue,urlcolor=pmblue,citecolor=pmblue,
  pdftitle={ProtMate User Guide},pdfauthor={Sanjeeva Wijeyesakere}}

\pagestyle{fancy}
\fancyhf{}
\fancyhead[L]{\small ProtMate \pmversion{} User Guide}
\fancyhead[R]{\small\leftmark}
\fancyfoot[C]{\thepage}
\renewcommand{\headrulewidth}{0.4pt}
\fancypagestyle{plain}{\fancyhf{}\fancyfoot[C]{\thepage}\renewcommand{\headrulewidth}{0pt}}

\setlist{itemsep=2pt,topsep=4pt}
\setlength{\parindent}{0pt}
\setlength{\parskip}{6pt}
\setlength{\emergencystretch}{3em}
\renewcommand{\arraystretch}{1.25}
\captionsetup{font=small,labelfont=bf}

% Text helpers
\newcommand{\cmd}[1]{\texttt{#1}}            % something to type, or a file name
\newcommand{\key}[1]{\fbox{\footnotesize\textsf{#1}}}
\newcommand{\gui}[1]{\textsf{\textbf{#1}}}  % a button, field or menu in the GUI
\newcommand{\menupath}[2]{\gui{#1}\,$\triangleright$\,\gui{#2}}
\newcommand{\ul}{\ensuremath{\mu}l}
\newcommand{\degC}{\ensuremath{^{\circ}}C}

% Note, tip and warning boxes
\newtcolorbox{note}[1][Note]{colback=pmgrey,colframe=pmblue,boxrule=0.5pt,arc=2pt,
  left=6pt,right=6pt,top=4pt,bottom=4pt,fonttitle=\bfseries\small,title=#1}
\newtcolorbox{warning}[1][Caution]{colback=pmwarn!6,colframe=pmwarn,boxrule=0.5pt,arc=2pt,
  left=6pt,right=6pt,top=4pt,bottom=4pt,fonttitle=\bfseries\small,title=#1}

% Terminal sessions and listings. In a session, \inp{...} marks what the user types.
\newcommand{\inp}[1]{\textbf{\textcolor{pmblue}{#1}}}
\DefineVerbatimEnvironment{session}{Verbatim}{fontsize=\small,frame=single,
  framesep=3mm,rulecolor=\color{gray},commandchars=\\\{\}}
\DefineVerbatimEnvironment{terminal}{Verbatim}{fontsize=\small,frame=leftline,
  framesep=3mm,rulecolor=\color{pmblue}}
\DefineVerbatimEnvironment{listing}{Verbatim}{fontsize=\small,frame=single,
  framesep=3mm,rulecolor=\color{gray}}

% A screenshot
\newcommand{\screenshot}[2]{%
  \begin{figure}[htbp]\centering
  \IfFileExists{images/#1.png}{\includegraphics[width=0.92\linewidth]{images/#1.png}}%
    {\fbox{\parbox[c][4cm][c]{0.8\linewidth}{\centering Screenshot images/#1.png not found}}}
  \caption{#2}\label{fig:#1}
  \end{figure}}

% Prompt tables for the terminal modules
\newenvironment{prompts}
  {\tabularx{\linewidth}{@{}>{\raggedright\arraybackslash\ttfamily\small}p{0.46\linewidth}X@{}}
   \toprule \textrm{\textbf{Prompt}} & \textbf{What to enter}\\ \midrule}
  {\bottomrule\endtabularx}
% The command names that start a terminal tool
% Heading for one item of the terminal version's main menu (numbered by the menu itself)
\newcommand{\menuoption}[1]{\subsection*{#1}\addcontentsline{toc}{subsection}{#1}}
\newcommand{\typed}[1]{\par\noindent\textbf{Type:} #1\par}

\title{%
  \vspace*{2cm}
  \IfFileExists{../resources/pmate_icon.png}{\includegraphics[width=3.5cm]{../resources/pmate_icon.png}\\[1cm]}{}
  {\Huge\bfseries ProtMate}\\[0.4cm]
  {\LARGE Laboratory Formula Calculator}\\[0.8cm]
  {\Large User Guide for version \pmversion}\\[0.3cm]
  {\large macOS \textperiodcentered{} Windows \textperiodcentered{} Linux}}
\author{Sanjeeva Wijeyesakere}
\date{October 2026}

\begin{document}

\maketitle
\thispagestyle{empty}

\clearpage
\thispagestyle{empty}
\vspace*{\fill}
{\small
\textbf{ProtMate \pmversion{} User Guide}

Copyright \textcopyright{} \copyrightyears{} Sanjeeva Wijeyesakere

ProtMate is free software: you can redistribute it and/or modify it under the terms of the GNU
General Public License as published by the Free Software Foundation, \gplgrant.

ProtMate is distributed in the hope that it will be useful, but WITHOUT ANY WARRANTY; without even
the implied warranty of MERCHANTABILITY or FITNESS FOR A PARTICULAR PURPOSE. See the GNU General
Public License for more details. You should have received a copy of the GNU General Public License
along with ProtMate. If not, see \url{https://www.gnu.org/licenses/}.

\textbf{Disclaimer.} ProtMate and this guide are provided ``as is''. The author assumes no
liability for the accuracy of any calculation, result or report produced by ProtMate, or for the
functionality, fitness or continued availability of the program. You are responsible for checking
all results independently before relying on them in experiments, publications or any other
decisions. See Appendix~\ref{app:licence} for the full terms.

Source code, releases and issue tracker: \repo
}
\clearpage

\tableofcontents

% ===========================================================================
\chapter{Introduction}

\section{What ProtMate does}

ProtMate is a calculator for the everyday arithmetic of a biochemistry or structural biology
laboratory. It works out protein and nucleic acid concentrations from absorbance readings, analyses
protein sequences (composition, chemical formula, molecular weight and extinction coefficient),
estimates primer melting temperatures and converts between percentage and molar concentrations. It
also includes small utilities for crystallography and docking work: renaming diffraction images for
Mosflm, writing AutoDock Vina configuration files, converting PDBQT files to PDB, and generating
X-PLOR/CNS scripts that solvate a protein.

\section{Two ways to use ProtMate}

ProtMate comes in two forms that share the same calculation engine, so they always give the same
answers and write identical report files.

\begin{description}
  \item[ProtMate (graphical version).] A desktop application with a window, a list of tools on the
    left and forms on the right. Results update as you type. This is the easiest way to start.
  \item[pmate (terminal version).] A text program that you run in a terminal window
    (Terminal on macOS, Command Prompt or PowerShell on Windows, a shell on Linux). You choose a tool
    from a numbered menu and answer its questions one at a time. It is handy on servers, over
    SSH connections, and for people who prefer the keyboard.
\end{description}

\begin{table}[htbp]\centering\small
\caption{The graphical and terminal versions compared.}
\begin{tabularx}{\linewidth}{@{}lXX@{}}
\toprule
 & \textbf{ProtMate (GUI)} & \textbf{pmate (terminal)} \\ \midrule
Program file & \cmd{ProtMate.app} (macOS), \cmd{ProtMate.exe} (Windows), \cmd{ProtMate} (Linux)
             & \cmd{pmate} (macOS, Linux), \cmd{pmate.exe} (Windows) \\
How you work & Forms; results update live & Question and answer, one value at a time \\
Calculators & All & All \\
File utilities & All, with previews before anything is written & All \\
Operating-system commands (\cmd{ls}, \cmd{ps}, \cmd{sys}, \ldots) & Not available & Available \\
Needs & A desktop session & Only a terminal \\
\bottomrule
\end{tabularx}
\end{table}

\section{Features at a glance}

Table~\ref{tab:features} lists every tool, the number you type for it in the terminal version's
menu, and the page that offers it in the graphical version.

\begin{table}[htbp]\centering\small
\caption{ProtMate's tools.}\label{tab:features}
\begin{tabularx}{\linewidth}{@{}Xcl@{}}
\toprule
\textbf{Tool} & \textbf{Terminal menu} & \textbf{GUI page} \\ \midrule
Mass needed for a \% (w/v) solution & 1 & Weight to Volume Ratio \\
Volume needed for a \% (w/v) solution & 2 & Weight to Volume Ratio \\
Protein concentration from one A280 reading & 3 & Protein Concentration (A280) \\
Protein concentration from several A280 readings (typed) & 4 & Multiple A280 Readings \\
Protein concentration from several A280 readings (from a file) & 5 & Multiple A280 Readings (\gui{Import\ldots}) \\
Extinction coefficient from residue counts & 6 & Extinction Coefficient \\
Full analysis of a protein sequence & 7 & Sequence Analysis \\
Nucleic acid concentration from A260 & 8 & Nucleic Acid Concentration (A260) \\
Primer/oligonucleotide melting temperature & 9 & Primer Melting Point \\
Body mass index & 10 & Body Mass Index \\
Factorial & 11 & Factorial \\
Mass/volume percent to molar concentration & \cmd{pmol}\textsuperscript{*} & Mass/Volume Percent to Molar Conc. \\ \midrule
Rename oscillation (diffraction) images for Mosflm & 12 & Rename Oscillation Files \\
AutoDock Vina configuration file (and run Vina) & 13 & AutoDock Vina Configuration \\
PDBQT to PDB conversion & 14 & PDBQT to PDB \\
X-PLOR/CNS code to solvate a protein & 15 & Solvate a Protein (X-PLOR) \\
Open a structure in RasMol & \cmd{open} & \menupath{File}{Open in RasMol\ldots} \\
\bottomrule
\multicolumn{3}{@{}l}{\footnotesize\textsuperscript{*} Typed at the prompt; not shown in the numbered menu.}
\end{tabularx}
\end{table}

\section{Licence and disclaimer}

ProtMate is free and open-source software, released under the GNU General Public License
version~3 (GPLv3). You may use it for any purpose, study and change its source code, and share
copies, original or modified, provided that you pass on the same freedoms under the same licence.
The source code is published at \repo.

\begin{warning}[Disclaimer]
ProtMate is provided without any warranty. The author assumes no liability for the accuracy of its
calculations or for the functionality of the program. Always check results that matter (for example
concentrations used to prepare samples, or values you intend to publish) by an independent method.
The body mass index calculator is included for convenience only and is not medical advice.
Appendix~\ref{app:licence} gives the full terms.
\end{warning}

\section{Conventions in this guide}

\begin{itemize}
  \item \cmd{Monospaced text} is something you type, or the name of a file or program.
  \item \gui{Bold sans-serif text} is a button, field or menu in the graphical version, and
        \menupath{File}{Quit} means ``open the \gui{File} menu and choose \gui{Quit}''.
  \item \key{Return} is a key on the keyboard. On some keyboards it is labelled \key{Enter}.
  \item In example terminal sessions, \inp{blue bold text} is what you type; everything else is
        printed by ProtMate. The examples use the user name \cmd{alice}; you will see your own.
\end{itemize}

% ===========================================================================
\chapter{Getting ProtMate}

\section{System requirements}

\begin{table}[htbp]\centering\small
\caption{Requirements for the ready-to-run versions.}\label{tab:requirements}
\begin{tabularx}{\linewidth}{@{}lX@{}}
\toprule
\textbf{Platform} & \textbf{Requirements} \\ \midrule
macOS & A Mac with Apple silicon (M1 or later) running macOS 27 or later. About 60\,MB of
        disk space. For an Intel Mac or an earlier macOS version, build ProtMate from source
        (Chapter~\ref{ch:source}). \\
Windows & 64-bit Windows 10 or Windows 11 on an x64 (Intel or AMD) processor. About 35\,MB of
          disk space. Nothing needs to be installed and administrator rights are not required. \\
Linux & No ready-made package; build from source (Chapter~\ref{ch:linux}). You need a C++17
        compiler, CMake 3.21 or later and, for the graphical version, Qt~6.5 or later. \\
\bottomrule
\end{tabularx}
\end{table}

AutoDock Vina and RasMol are optional. ProtMate works without them; they are only needed for the
``run Vina'' and ``open in RasMol'' features (Chapter~\ref{ch:external}).

\section{What to download}

ProtMate is published on GitHub at \repo. Ready-to-run versions are attached to each release on
the \gui{Releases} page, \releases. Open the latest release and, under \gui{Assets}, download the
file for your computer:

\begin{table}[htbp]\centering\small
\caption{ProtMate distribution files.}
\begin{tabularx}{\linewidth}{@{}lX@{}}
\toprule
\textbf{File} & \textbf{Contents} \\ \midrule
\cmd{ProtMate-\pmversion.dmg} & macOS disk image containing the graphical version
                                (\cmd{ProtMate.app}). \\
\cmd{pmate} & macOS terminal version (a single file). \\
\cmd{ProtMate-\pmversion-win64-portable.zip} & Windows portable package with both the graphical
                                (\cmd{ProtMate.exe}) and terminal (\cmd{pmate.exe}) versions. \\
\gui{Source code} (zip or tar.gz) & Everything needed to build ProtMate yourself on any
                                platform: \cmd{core/}, \cmd{cli/}, \cmd{gui/}, \cmd{resources/},
                                \cmd{tests/}, \cmd{docs/} and \cmd{CMakeLists.txt}. GitHub adds
                                these archives to every release automatically. \\
\bottomrule
\end{tabularx}
\end{table}

To get the latest, unreleased source code instead, use \gui{Code}\,$\triangleright$\,\gui{Download
ZIP} on the repository page, or clone the repository with Git:

\begin{terminal}
git clone https://github.com/swijeyes/ProtMate.git
\end{terminal}

% ===========================================================================
\chapter{Installing on macOS}\label{ch:mac}

\section{Installing the graphical version}

\begin{enumerate}
  \item Double-click \cmd{ProtMate-\pmversion.dmg} in your Downloads folder. A window opens
        showing the ProtMate icon and a shortcut to your Applications folder.
  \item Drag the \gui{ProtMate} icon onto the \gui{Applications} shortcut.
  \item Eject the disk image: click the eject symbol next to \gui{ProtMate \pmversion} in the
        Finder sidebar, or drag the disk icon on your desktop to the Trash.
  \item Open ProtMate from Launchpad, Spotlight (\key{Cmd}+\key{Space}, then type
        \cmd{ProtMate}), or the Applications folder.
\end{enumerate}

\section{The first time you open ProtMate}\label{sec:gatekeeper}

ProtMate is not signed with an Apple Developer ID, so the first time you open it macOS says it
cannot verify the developer and offers only \gui{Done} or \gui{Move to Trash}. To allow it:

\begin{enumerate}
  \item Click \gui{Done} (not \gui{Move to Trash}).
  \item Open \menupath{System Settings}{Privacy \& Security} and scroll down to the
        \gui{Security} section. You will see a message that ProtMate was blocked.
  \item Click \gui{Open Anyway}, confirm with your password or Touch ID, and click
        \gui{Open Anyway} again in the dialog that follows.
\end{enumerate}

You only need to do this once. Alternatively, an administrator can clear the download flag in
Terminal:

\begin{terminal}
xattr -dr com.apple.quarantine /Applications/ProtMate.app
\end{terminal}

\section{Installing the terminal version}

The terminal version is the single file \cmd{pmate}. Put it in a folder on your command search
path so you can run it from anywhere. Open Terminal and run:

\begin{terminal}
mkdir -p ~/bin
mv ~/Downloads/pmate ~/bin/
chmod +x ~/bin/pmate
xattr -d com.apple.quarantine ~/bin/pmate   # clears the download flag
\end{terminal}

If \cmd{\textasciitilde/bin} is not already on your search path, add it (the default shell on
macOS is zsh) and open a new Terminal window:

\begin{terminal}
echo 'export PATH="$HOME/bin:$PATH"' >> ~/.zshrc
\end{terminal}

Now type \cmd{pmate} in any Terminal window to start the terminal version. To install it for all
users instead, copy it to \cmd{/usr/local/bin} with \cmd{sudo cp pmate /usr/local/bin/}.

\section{Uninstalling}

Drag \cmd{ProtMate} from the Applications folder to the Trash and delete \cmd{pmate} from wherever
you put it. The only other thing ProtMate keeps is a small preferences file that remembers the last
folder you used.

% ===========================================================================
\chapter{Installing on Windows}\label{ch:windows}

The Windows version is \emph{portable}: it runs from its own folder and makes no changes to the
rest of the system, so it does not need an installer or administrator rights.

\section{Unpacking}

\begin{enumerate}
  \item Right-click \cmd{ProtMate-\pmversion-win64-portable.zip} in your Downloads folder and
        choose \gui{Extract All\ldots}.
  \item Choose where to keep ProtMate, for example \cmd{C:\textbackslash Users\textbackslash
        \emph{you}\textbackslash Apps}, and click \gui{Extract}. This creates a folder called
        \cmd{ProtMate-\pmversion-win64}.
\end{enumerate}

Do not run ProtMate from inside the zip file without extracting it first, and keep the folder's
contents together:

\begin{listing}
ProtMate-1.2.0-win64\
    ProtMate.exe           graphical version
    pmate.exe              terminal version
    Qt6Core.dll, Qt6Gui.dll, Qt6Widgets.dll, libc++.dll, libunwind.dll
    platforms\qwindows.dll
    styles\qmodernwindowsstyle.dll
    licenses\              licence texts for the bundled libraries
    README.txt
    LICENSE.txt            the GNU General Public License (ProtMate's licence)
\end{listing}

\cmd{ProtMate.exe} needs every file in the folder. \cmd{pmate.exe} is self-contained and can be
copied elsewhere on its own.

\section{The first time you start ProtMate}

Double-click \cmd{ProtMate.exe}. Because the program is not code-signed, Microsoft Defender
SmartScreen may show ``Windows protected your PC''. Click \gui{More info} and then
\gui{Run anyway}. Windows remembers your choice.

To make ProtMate easier to find, right-click \cmd{ProtMate.exe} and choose \gui{Pin to Start}, or
\gui{Show more options}\,$\triangleright$\,\gui{Send to}\,$\triangleright$\,\gui{Desktop (create shortcut)}.

\section{Using the terminal version on Windows}

\cmd{pmate.exe} is a console program. You can double-click it, which opens it in a console window,
but it is usually better to start it from a terminal so the window stays open after you quit:

\begin{enumerate}
  \item Open the ProtMate folder in File Explorer, click in the address bar, type \cmd{cmd} and
        press \key{Enter}. A Command Prompt opens in that folder. (In Windows~11 you can also
        right-click inside the folder and choose \gui{Open in Terminal}.)
  \item Type \cmd{pmate} and press \key{Enter}.
\end{enumerate}

To run \cmd{pmate} from any folder, add the ProtMate folder to your user \cmd{Path}: open the
Start menu, search for \gui{Edit environment variables for your account}, select \gui{Path},
click \gui{Edit\ldots}, then \gui{New}, and paste the folder's location. Open a new terminal
window afterwards.

\begin{note}
The terminal version saves files and looks for input files in the \emph{current folder} of the
terminal. Use \cmd{cd} to move to the folder that holds your data before starting \cmd{pmate}.
\end{note}

\section{Uninstalling}

Delete the ProtMate folder and any shortcuts you made. ProtMate also keeps the last folder you
used in the registry key \cmd{HKEY\_CURRENT\_USER\textbackslash Software\textbackslash
ProtMate}, which you can delete with the Registry Editor if you wish.

% ===========================================================================
\chapter{Installing on Linux}\label{ch:linux}

There is no ready-made Linux package, so you build ProtMate from its source code. This takes a few
minutes and produces both the graphical (\cmd{ProtMate}) and the terminal (\cmd{pmate}) versions.

\section{Install the build tools}

The graphical version needs Qt~6.5 or later. Recent distributions provide it, for example Debian~13,
Ubuntu~24.10 and later, Fedora~39 and later, and Arch Linux. Debian~12 and Ubuntu~22.04/24.04 ship
an older Qt~6, so on those either build only the terminal version
(Section~\ref{sec:cli-only}) or install a newer Qt from \url{https://www.qt.io/download-open-source}.

\begin{description}
  \item[Debian, Ubuntu and derivatives]\leavevmode
\begin{terminal}
sudo apt update
sudo apt install build-essential cmake qt6-base-dev
\end{terminal}
  \item[Fedora]\leavevmode
\begin{terminal}
sudo dnf install gcc-c++ cmake qt6-qtbase-devel
\end{terminal}
  \item[Arch Linux and derivatives]\leavevmode
\begin{terminal}
sudo pacman -S --needed base-devel cmake qt6-base
\end{terminal}
\end{description}

To check the Qt version, run \cmd{qmake6 -{}-version} or look for the package version in your
package manager.

\section{Build}

Download the source code from the latest release (\releases) and unpack it, or clone the
repository (this needs \cmd{git}, which you can install with your package manager):

\begin{terminal}
git clone https://github.com/swijeyes/ProtMate.git
cd ProtMate
\end{terminal}

Then, in a terminal in the folder that contains \cmd{CMakeLists.txt}, run:

\begin{terminal}
cmake -S . -B build
cmake --build build -j
ctest --test-dir build        # optional: runs ProtMate's self-tests
\end{terminal}

The last command should finish with \cmd{100\% tests passed}. The programs are now in the
\cmd{build} folder; try them with \cmd{./build/ProtMate} and \cmd{./build/pmate}.

\section{Install}

To install both programs into \cmd{/usr/local/bin}:

\begin{terminal}
sudo cmake --install build
\end{terminal}

Use \cmd{cmake -{}-install build -{}-prefix \textasciitilde/.local} instead to install just for
yourself without \cmd{sudo} (the programs go to \cmd{\textasciitilde/.local/bin}, which most
distributions already have on the search path).

\section{Adding ProtMate to the applications menu}

To start the graphical version from your desktop's application launcher, copy the icon and create
a desktop entry:

\begin{terminal}
mkdir -p ~/.local/share/icons ~/.local/share/applications
cp resources/pmate_icon.png ~/.local/share/icons/protmate.png
cat > ~/.local/share/applications/protmate.desktop <<'EOF'
[Desktop Entry]
Type=Application
Name=ProtMate
Comment=Laboratory formula calculator
Exec=ProtMate
Icon=protmate
Categories=Science;Education;
EOF
\end{terminal}

If you installed with \cmd{-{}-prefix \textasciitilde/.local}, change the \cmd{Exec} line to the
full path, for example \cmd{Exec=/home/\emph{you}/.local/bin/ProtMate}.

\section{Building only the terminal version}\label{sec:cli-only}

The terminal version needs no Qt at all. This works on any Linux system with a C++17 compiler and
CMake, including servers without a desktop:

\begin{terminal}
cmake -S . -B build -DPMATE_BUILD_GUI=OFF
cmake --build build -j
sudo cmake --install build
\end{terminal}

\section{Uninstalling}

CMake records what it installed in \cmd{build/install\_manifest.txt}. To remove it:

\begin{terminal}
xargs sudo rm -v < build/install_manifest.txt
\end{terminal}

Also delete these files if you created or used them:
\begin{itemize}
  \item \cmd{\textasciitilde/.local/share/applications/protmate.desktop}
  \item \cmd{\textasciitilde/.local/share/icons/protmate.png}
  \item \cmd{\textasciitilde/.config/ProtMate/ProtMate.conf} (ProtMate's settings)
\end{itemize}

% ===========================================================================
\chapter{Building from source on any platform}\label{ch:source}

The same steps work on macOS and Windows, which is useful on an Intel Mac, an older macOS version,
or to make a build of your own. Get the source code from \repo{} as described in
Chapter~\ref{ch:linux}, and run the commands below in the folder that contains
\cmd{CMakeLists.txt}. Because ProtMate is licensed under the GPL, you are free to modify it; if you
distribute a modified version, you must also make its source code available under the GPL.

\section{Requirements}

\begin{itemize}
  \item CMake 3.21 or later.
  \item A C++17 compiler: Apple Clang (Xcode or the Command Line Tools) on macOS, GCC or Clang on
        Linux, Microsoft Visual C++ 2019 or later, or MinGW on Windows.
  \item Qt~6.5 or later with the Widgets module, for the graphical version.
\end{itemize}

\section{macOS}

With Homebrew (\url{https://brew.sh}):

\begin{terminal}
xcode-select --install          # if you do not already have a compiler
brew install cmake qt
cmake -S . -B build -DCMAKE_PREFIX_PATH="$(brew --prefix qt)"
cmake --build build -j
open build/ProtMate.app
\end{terminal}

To make a self-contained copy of the app that you can move to another Mac, run
\cmd{macdeployqt build/ProtMate.app}. The \cmd{README.md} file in the source folder describes the
remaining packaging steps.

\section{Windows}

Install Qt with the Qt Online Installer (\url{https://www.qt.io/download-open-source}), choosing
the MSVC or MinGW build that matches your compiler, and install CMake. Then, in a developer command
prompt:

\begin{terminal}
cmake -S . -B build -DCMAKE_PREFIX_PATH=C:\Qt\6.11.2\msvc2022_64
cmake --build build --config Release
C:\Qt\6.11.2\msvc2022_64\bin\windeployqt build\Release\ProtMate.exe
\end{terminal}

Adjust the Qt path to the version and compiler you installed.

\section{Build options}

\begin{table}[htbp]\centering\small
\caption{Options for the \cmd{cmake -S . -B build} step.}
\begin{tabularx}{\linewidth}{@{}lX@{}}
\toprule
\textbf{Option} & \textbf{Effect} \\ \midrule
\cmd{-DPMATE\_BUILD\_GUI=OFF} & Build only the terminal version; Qt is not needed. \\
\cmd{-DPMATE\_BUILD\_TESTS=OFF} & Skip the self-tests. \\
\cmd{-DPMATE\_SANITIZE=ON} & Developer builds with AddressSanitizer and UndefinedBehaviorSanitizer. \\
\cmd{-DCMAKE\_PREFIX\_PATH=\emph{dir}} & Where to find Qt, if CMake cannot find it by itself. \\
\cmd{-DCMAKE\_BUILD\_TYPE=Debug} & Debugging build (the default is \cmd{Release}). \\
\bottomrule
\end{tabularx}
\end{table}

% ===========================================================================
\chapter{Optional programs: AutoDock Vina and RasMol}\label{ch:external}

ProtMate can start two external programs if they are installed:

\begin{description}
  \item[AutoDock Vina] (\url{https://vina.scripps.edu}) runs a docking calculation with the
        configuration file ProtMate writes.
  \item[RasMol] (\url{http://www.openrasmol.org}) displays a PDB structure.
\end{description}

ProtMate looks for them on your command search path (\cmd{PATH}) under the names \cmd{vina} and
\cmd{rasmol} (\cmd{vina.exe} and \cmd{rasmol.exe} on Windows). Only folders given as full paths
are searched, never the current folder, so a stray program in your data folder cannot be started
by mistake. To check that a program can be found, run \cmd{which vina} (macOS, Linux) or
\cmd{where vina} (Windows) in a terminal. If nothing is printed, add the program's folder to
\cmd{PATH} and restart ProtMate.

When a program cannot be found, the graphical version disables the corresponding button or menu
item and says so, and the terminal version prints a message. Everything else keeps working; for
example you can still write a Vina configuration file and run Vina yourself.

\begin{note}
On macOS, applications started from the Dock or Finder do not see changes you make to
\cmd{PATH} in \cmd{\textasciitilde/.zshrc}. If the graphical version does not find Vina or RasMol
although \cmd{which} does, start it from Terminal with \cmd{open -a ProtMate}, or install the
program in \cmd{/usr/local/bin}.
\end{note}

% ===========================================================================
\chapter{Quick start}

\section{Graphical version}

\begin{enumerate}
  \item Start ProtMate. The window shows a list of tools on the left.
  \item Click \gui{Protein Concentration (A280)}.
  \item Fill in the fields: total volume of protein \cmd{2.5}\,ml, A280 \cmd{0.42}, total liquid in
        the cuvette \cmd{100}\,\ul, protein in the cuvette \cmd{10}\,\ul{} and extinction
        coefficient \cmd{0.92}.
  \item The results appear immediately: dilution factor 10, protein concentration 4.57\,mg/ml and
        11.425\,mg of protein in total.
\end{enumerate}

\section{Terminal version}

\begin{enumerate}
  \item Open a terminal and type \cmd{pmate}, then press \key{Return}.
  \item Type \cmd{mnu} and press \key{Return} to see the main menu.
  \item Type \cmd{3} and press \key{Return}, then answer the questions with the numbers above.
  \item Read the results, press \key{Return} to go back, and type \cmd{0} to quit.
\end{enumerate}

\begin{session}
alice% > \inp{3}

Enter total volume of protein (in ml): \inp{2.5}
Enter A280 of protein: \inp{0.42}
Enter total volume of liquid in cuvette (protien + buffer) (in ul): \inp{100}
Enter volume of protein in cuvette (ul): \inp{10}
Enter protein's extinction coefficient: \inp{0.92}

Dilution factor: 10

Protein concentration: 4.57 mg/ml
Total amount of protein present in 2.5 ml: 11.425 mg

Press the Return key to continue....
\end{session}

% ===========================================================================
\chapter{Using the graphical version}\label{ch:gui}

\section{The main window}

\screenshot{gui_main}{The ProtMate window. Tools are listed on the left, grouped into
\gui{Calculators} and \gui{File Utilities}; the selected tool fills the right-hand side.}

The list on the left contains every tool, in the same groups as the terminal version's menu. Click
a tool to open it; the values you entered on each page are kept while ProtMate is running, so you
can move between tools freely. You can drag the divider between the list and the page to change
their widths, and pages scroll if the window is too small for them.

\section{How the pages work}

\begin{itemize}
  \item \textbf{Results update as you type.} There is no ``Calculate'' button.
  \item \textbf{Number fields} accept only sensible values and show their units. Type a value, or use
        the arrows, the \key{$\uparrow$}/\key{$\downarrow$} keys or the mouse wheel to change it.
  \item \textbf{Problems are explained on the page}, for example an unknown residue in a sequence
        or a file name that is not allowed, and the save buttons stay disabled until they are
        fixed.
  \item \textbf{Save report\ldots} buttons open the standard save dialog. The file has exactly the
        same contents as the report the terminal version writes. The save dialog asks before
        replacing an existing file, and a file is only replaced once the new version has been
        written completely.
  \item \textbf{Open buttons} (\gui{Open .seq file\ldots}, \gui{Browse\ldots}, \gui{Import\ldots})
        use the standard open dialog, which starts in the folder you used last.
\end{itemize}

\section{Menus}

\begin{description}
  \item[\menupath{File}{Open in RasMol\ldots}] Choose a PDB file to display in RasMol. The item is
        disabled, and says ``RasMol not found'', if RasMol is not installed
        (Chapter~\ref{ch:external}).
  \item[\menupath{File}{Quit}] Close ProtMate (\key{Cmd}+\key{Q} on macOS, \key{Ctrl}+\key{Q}
        elsewhere). If AutoDock Vina is running, it is stopped.
  \item[\menupath{Help}{About ProtMate}] Shows the version, the copyright and licence notice, the
        disclaimer, a link to the project on GitHub, and information about your system.
  \item[\menupath{Help}{About Qt}] Shows the version and licence of the Qt toolkit that the
        graphical version is built with.
\end{description}
On macOS, both \gui{About} items are in the \gui{ProtMate} application menu.

\section{Calculators}

\subsection{Weight to Volume Ratio}

Two small calculators for \% (w/v) solutions (see Figure~\ref{fig:gui_main}):

\begin{itemize}
  \item \gui{Calculate weight}: enter the weight to volume percent and the desired volume (ml) to
        get the mass of solid to weigh out (g). For example, 5\% in 200\,ml needs 10\,g.
  \item \gui{Calculate volume}: enter the percent and the weight of sample you have (g) to get the
        volume of liquid to dissolve it in (ml).
\end{itemize}

\subsection{Protein Concentration (A280)}

\screenshot{gui_protein}{Protein concentration from a single A280 reading.}

Calculates the concentration and total amount of a protein from one absorbance reading at 280\,nm.

\begin{description}
  \item[Total volume of protein] The volume of your whole protein sample, in ml.
  \item[A280 of protein] The absorbance reading.
  \item[Total liquid in cuvette] The total volume in the cuvette (protein plus buffer), in \ul.
  \item[Protein in cuvette] How much of your protein sample is in the cuvette, in \ul. Together
        with the total, this gives the dilution factor.
  \item[Extinction coefficient] The absorbance at 280\,nm of a 1\,mg/ml solution of your protein
        in a 1\,cm cuvette (sometimes written $A_{280}^{0.1\%}$). The \gui{Extinction Coefficient}
        and \gui{Sequence Analysis} pages calculate this value for you; it is the figure they
        label ``Absorbance (280\,nm, 1\,mg/ml, 1\,cm)''.
\end{description}

The results are the dilution factor, the protein concentration (mg/ml) and the total amount of
protein (mg). See Appendix~\ref{app:formulas} for the formulas.

\subsection{Multiple A280 Readings}

\screenshot{gui_multi}{Several readings at once, for example the fractions from a column.}

Use this page when you have several samples measured the same way, such as fractions from a
purification. Enter the cuvette volumes and extinction coefficient once at the top, then one row
per sample with its A280 and its volume (ml). The concentration and yield of each sample are
filled in as you type.

\begin{itemize}
  \item \gui{Add reading} adds a row; \gui{Remove selected} deletes the selected rows;
        \gui{Clear} removes them all.
  \item \gui{Import\ldots} reads readings from a text file (Section~\ref{sec:readings-file}),
        replacing the rows in the table. Up to 100{,}000 readings are read.
  \item \gui{Save report\ldots} saves a tab-separated table that opens directly in a spreadsheet.
  \item Rows that are incomplete or contain something other than a number are reported below the
        table, and must be fixed before saving.
\end{itemize}

\subsection{Extinction Coefficient}

\screenshot{gui_extinction}{Extinction coefficient from residue counts.}

Enter the number of tyrosines, tryptophans and cysteines and the molecular weight (Da). The page
shows the molar extinction coefficient of the native protein in water
(M$^{-1}$\,cm$^{-1}$) and the absorbance at 280\,nm of a 1\,mg/ml solution in a 1\,cm cuvette.
\gui{Save report\ldots} saves both values. If you have the protein's sequence, the
\gui{Sequence Analysis} page counts the residues for you.

\subsection{Sequence Analysis}

\screenshot{gui_sequence}{Analysing a protein sequence.}

Analyses a protein or peptide sequence written in one-letter amino acid codes.

\begin{enumerate}
  \item Paste the sequence into the box, or click \gui{Open .seq file\ldots} to load one
        (Section~\ref{sec:seq-files}). Spaces and line breaks are ignored, and upper or lower case
        both work.
  \item \gui{Residues in sequence} shows the length. To analyse part of the sequence (a domain, for
        example), change \gui{First residue} and \gui{Final residue}; they start at the whole
        sequence.
  \item The report below shows the amino acid composition (number and percentage of each residue),
        the chemical formula, the molecular weight, the extinction coefficient and the absorbance of
        a 1\,mg/ml solution.
  \item To save the report, type a name for the protein in the box at the bottom (it appears in the
        report's heading) and click \gui{Save report\ldots}. The report also records your user
        name as the operator.
\end{enumerate}

Characters that are not one of the 20 standard amino acid codes (for example \cmd{X}, \cmd{B},
\cmd{Z}, \cmd{*} or digits) are counted as \cmd{Unk}. Their positions are listed at the top of
the report with a warning to check the sequence, and they contribute nothing to the formula or
molecular weight.

\subsection{Nucleic Acid Concentration (A260)}

\screenshot{gui_nucleic}{Nucleic acid concentration from A260.}

Choose dsDNA, RNA or Oligonucleotide, enter the OD at 260\,nm and the dilution factor (a whole
number; 1 if undiluted). The concentration is shown in \textmu g/ml. A warning appears when the
reading is outside 0.1--10, where absorbance readings become unreliable.

\subsection{Primer Melting Point}

\screenshot{gui_melt}{Primer melting temperature.}

Estimates the melting temperature ($T_m$) of a primer or oligonucleotide. Choose how to describe it:

\begin{itemize}
  \item \gui{Enter length and G+C count}: give the number of nucleotides and how many are G or C.
  \item \gui{From a sequence}: paste a DNA sequence or click \gui{Open .seq file\ldots}. Only A, C,
        G and T are allowed (in either case); the first other character is reported with its
        position. The counts of each base are shown with the results.
\end{itemize}

Enter the total salt (Na$^+$) concentration in molar; 0.05\,M is typical. The page shows the G--C
content, the basic melting point and the salt-adjusted melting point. \gui{Save report\ldots}
saves them.

\subsection{Mass/Volume Percent to Molar Conc.}

\screenshot{gui_percent}{Converting a percentage to a molar concentration.}

Converts a percentage solution to molarity. Choose:

\begin{itemize}
  \item \gui{Mass concentration (\% w/v)}: enter the percentage and the formula weight (g/mol).
        For example, 5\% (w/v) NaCl (58.44\,g/mol) is 0.86\,M (855.58\,mM).
  \item \gui{Volume percent (\% v/v)}, for liquids: also enter the density of the pure compound
        (g/ml).
\end{itemize}

The result is shown in both M and mM.

\subsection{Body Mass Index}

Enter height in feet and inches and weight in pounds. The page shows the height in metres, the
weight in kilograms, the body mass index with its category, and the weight range (lb) that
corresponds to a normal BMI for that height. The reference ranges from the US Centers for Disease
Control and Prevention are shown on the page.

\subsection{Factorial}

Enter a whole number from 0 to 20 to see its factorial. (21! is too large to be represented
exactly.)

\section{File utilities}

\subsection{Rename Oscillation Files}

\screenshot{gui_rename}{Renaming diffraction images, with a preview of every change.}

Renames a series of diffraction (oscillation) images into the \cmd{prefix\_\#\#\#.osc} form that
Mosflm reads. For example, \cmd{1.0001}, \cmd{1.0002}, \ldots{} become \cmd{a\_001.osc},
\cmd{a\_002.osc}, \ldots

\begin{enumerate}
  \item \gui{Directory}: click \gui{Browse\ldots} and choose the folder that holds the images.
  \item \gui{Number of files} and \gui{Initial frame number}: how many images there are and the
        frame number of the first one.
  \item \gui{Current file string}: the part of the current names before the frame number. For
        images called \cmd{1.0001}, \cmd{1.0002}, \ldots{} enter \cmd{1.}. ProtMate expects the
        current names to be this string followed by \cmd{0} and the frame number padded to three
        digits (\cmd{1.0001}).
  \item \gui{New file string}: the start of the new names; \cmd{a\_} gives \cmd{a\_001.osc}.
  \item Check the preview. Each file is marked \gui{Ready}, \gui{Not found}, or
        \gui{New name already exists}, and a summary appears below the table.
  \item Click \gui{Rename\ldots} and confirm. Only files marked \gui{Ready} are renamed.
\end{enumerate}

ProtMate never overwrites a file: if the new name already exists, that file is skipped. The file
strings may contain letters, digits and \cmd{. \_ - +}, and must not start with \cmd{-}.

\subsection{AutoDock Vina Configuration}

\screenshot{gui_vina}{Writing a configuration file for AutoDock Vina.}

Writes the configuration file that AutoDock Vina needs for a docking run, with a cubic search
box, and can then run Vina.

\begin{enumerate}
  \item Choose the \gui{Receptor} and \gui{Ligand} files (PDBQT format) with \gui{Browse\ldots}, or
        type their names. Names relative to a folder are interpreted from the folder where you save
        the configuration file.
  \item Enter the centre of the search box (\gui{Grid centre X}, \gui{Y}, \gui{Z}, in \AA) and the
        length of one side of the box (\AA).
  \item Check the \gui{Configuration preview}. Besides your values, the file sets the output
        file (\cmd{out = all.pdbqt}) and \cmd{exhaustiveness = 8}.
  \item Click \gui{Save configuration\ldots}. The suggested name is \cmd{conf.txt}, in the
        receptor's folder.
  \item If Vina is installed, click \gui{Run Vina}. Vina runs in the folder of the configuration
        file; its output appears in the log at the bottom of the page and is also written to
        \cmd{vina\_log.txt}. The docked poses are written to \cmd{all.pdbqt} in the same folder.
        \gui{Stop} ends a run early.
\end{enumerate}

\gui{Run Vina} is enabled only when the configuration shown has been saved, so the file Vina reads
always matches what you see. If you change a value, save again before running. File names must
not contain \cmd{\#} or control characters.

\subsection{PDBQT to PDB}

Converts a PDBQT file (the format AutoDock uses) to a standard PDB file in which the B-factor
column holds the partial atomic charges. Choose the PDBQT file, click
\gui{Convert and save\ldots}, and choose where to save the \cmd{.pdb} file.

\subsection{Solvate a Protein (X-PLOR)}

\screenshot{gui_xplor}{Generating X-PLOR/CNS scripts to solvate a protein.}

Generates X-PLOR (or CNS) input that overlays a protein on a pre-equilibrated solvent shell, in two
stages:

\begin{description}
  \item[Stage 1] reads the protein and solvent and writes \cmd{overlay1.pdb}, containing the
        solvent molecules within 6.0\,\AA{} of any protein heavy atom.
  \item[Stage 2] reads the protein and the solvent produced by stage 1, deletes solvent molecules
        within 2.3\,\AA{} of the protein (which would clash with it), and writes
        \cmd{overlay2.pdb}.
\end{description}

Choose \gui{Stage 1}, \gui{Stage 2} or \gui{Both stages} under \gui{Generate}, and fill in the
protein's PDB coordinate and MTF structure files and the solvent's files. For stage 2, the solvent
files are the ones produced by stage 1 (with \gui{Both stages}, they go in the
\gui{Stage 2 input} box). The scripts are shown on the \gui{Stage 1 script} and
\gui{Stage 2 script} tabs; save them with \gui{Save stage 1\ldots} and \gui{Save stage 2\ldots}
(\cmd{.inp} files) and run them in X-PLOR or CNS. The solvent segment is expected to have the
segment ID \cmd{W}.

\begin{note}
X-PLOR cannot handle spaces in file names, so the names in these scripts may contain only letters,
digits and \cmd{. \_ - / +}, and must not start with \cmd{-}. The page tells you if a name is not
allowed.
\end{note}

% ===========================================================================
\chapter{Using the terminal version}\label{ch:cli}

\section{Starting pmate}

Open a terminal, change to the folder that holds your data with \cmd{cd}, and type \cmd{pmate}:

\begin{description}
  \item[macOS] Open Terminal (in \menupath{Applications}{Utilities}), then for example
        \cmd{cd \textasciitilde/Documents/xtal} and \cmd{pmate}.
  \item[Linux] Open your terminal application, \cmd{cd} to your data and type \cmd{pmate}.
  \item[Windows] Open Command Prompt, PowerShell or Windows Terminal, \cmd{cd} to your data and type
        \cmd{pmate} (or the full path to \cmd{pmate.exe} if its folder is not on your \cmd{Path}).
\end{description}

The terminal version reads input files from, and saves files to, the current folder unless you type
a path. File names are typed without quotation marks.

\section{The prompt}

When \cmd{pmate} starts it shows a short heading and waits for a command:

\begin{session}
ProtMate
Developed by Sanjeeva Wijeyesakere
Ver 1.2.0       Session started on Sat Oct 25 09:15:02 PDT 2026

alice% >
\end{session}

The prompt \cmd{alice\% >} shows your user name. At the prompt you can type:

\begin{itemize}
  \item \cmd{mnu} (or \cmd{menu}) to show the main menu;
  \item a menu number, such as \cmd{3}, or a command name, such as \cmd{prot}, to start a tool;
  \item \cmd{?} or \cmd{help} for a list of commands;
  \item \cmd{0} (or \cmd{quit}, \cmd{exit}, \cmd{stop}, \cmd{bye}) to leave ProtMate.
\end{itemize}

Each command is followed by \key{Return}. Commands are case-sensitive: type them in lower case as
shown. Pressing \key{Return} on its own just shows the prompt again, and an unrecognised command
prints \cmd{Invalid Response ...}.

\section{The main menu}\label{sec:mnu}

Type \cmd{mnu} at the prompt to see the main menu:

\begin{session}
alice% > \inp{mnu}
ProtMate
Developed by Sanjeeva Wijeyesakere
Ver 1.2.0       Session started on Sat Oct 25 09:15:02 PDT 2026

ProtMate Main Menu
==================

Calculators:
 1.  Weight to Volume Ratio - Calculate weight
 2.  Weight to Volume Ratio - Calculate volume
 3.  Calculate Protein Concentration Using A280 (Single Reading)
 4.  Calculate Protein Concentration Using A280 (Multiple Readings)
 5.  Calculate Protein Concentration Using A280 (Multiple Readings, auto input)
 6.  Calculate Protein Extinction Coefficient
 7.  Calculate Protein Extinction Coefficient From a Sequence (*.seq) File
 8.  Calculate Concentration of a Nucleic Acid from A260
 9.  Calculate Melting Point of a Primer/Oligonucleotide
 10. Calculate your Body Mass Index (BMI)
 11. Factorial Calculator

External Interfaces & File Utilities:
 12. Rename Oscillation Files
 13. Generate Configuration/Input File to Run AutoDock Vina
 14. Convert a PDBQT file to a PDB file
 15. Generate XPLOR code to solvate a protein structure

Other Functions:
 ?   - Help (Command line commands)
 ls  - List Files in the Current Directory
 pwd - Display the name of the Current Directory
 sys - Issue a System Command
 abt - About Protmate
 0   - Quit

*** For a complete list of avaiable functions, consult the help menu
*** (enter "help" or "?" at the prompt)

alice% >
\end{session}

Type the number of a menu item and press \key{Return}. Each item also has a command name, which you
can type instead of the number (Table~\ref{tab:commands}). The menu is a reminder only: you can
type any number or command at the prompt at any time, whether or not the menu is on the screen.

\section{How the tools work}

\begin{itemize}
  \item A tool asks its questions one at a time. Type the answer and press \key{Return}.
  \item Numbers can be whole or decimal (\cmd{0.42}). Do not type units; each question says which
        unit it expects (for example ml or \ul).
  \item When a tool finishes, it usually shows \cmd{Press the Return key to continue....}. Press
        \key{Return} to clear the screen and go back to the prompt.
  \item Many tools let you type \cmd{0} at the first question to leave without doing anything.
        Tools that need a number at that point cannot offer this; you can always finish the
        questions and discard the result.
  \item Several tools offer to save their results. Answer \cmd{yes} (in full, in lower case) to save;
        anything else means no. Then type a file name. If the file already exists, ProtMate asks
        \cmd{Overwrite it (yes/no)?} and only replaces it if you type \cmd{yes}.
  \item Where a tool asks for a file name, type a single name or path without spaces, except where
        this guide says otherwise.
\end{itemize}

\begin{warning}
If you type something that is not a number where a number is expected, the value is taken as 0 and
the remaining questions of that tool may be skipped. The tool finishes and returns you to the
prompt, so just start it again. To quit ProtMate immediately at any time, press
\key{Ctrl}+\key{C}.
\end{warning}

% ---------------------------------------------------------------------------
\section{Menu options in detail}

The rest of this section describes each item of the main menu. The command names shown in brackets
can be typed instead of the number.

\menuoption{1. Weight to Volume Ratio -- Calculate weight}
\typed{\cmd{1}, \cmd{wv-weight}}

Works out how much solid to weigh out to make a solution of a given percentage (w/v).

\begin{prompts}
Enter weight to volume percent ( \%(w/v)): & The concentration you want, as a percentage, e.g. \cmd{5}. \\
Enter desired volume (in ml): & The final volume of the solution. \\
\end{prompts}

\begin{session}
alice% > \inp{1}

Enter weight to volume percent ( %(w/v)): \inp{5}
Enter desired volume (in ml): \inp{200}

Mass of sample required = 10 g
\end{session}

\menuoption{2. Weight to Volume Ratio -- Calculate volume}
\typed{\cmd{2}, \cmd{wv-vol}}

The reverse of option 1: how much liquid to dissolve a weighed sample in.

\begin{prompts}
Enter weight to volume percent ( \%(w/v)): & The concentration you want, as a percentage. \\
Enter sample weight (in g): & The mass of solid you have. \\
\end{prompts}

The result is printed as \cmd{Volume of liquid required = \ldots{} ml}.

\menuoption{3. Protein Concentration Using A280 (Single Reading)}
\typed{\cmd{3}, \cmd{prot}}

Calculates a protein's concentration and total amount from one absorbance reading at 280\,nm.

\begin{prompts}
Enter total volume of protein (in ml): & The volume of your whole protein sample. \\
Enter A280 of protein: & The absorbance reading. \\
Enter total volume of liquid in cuvette (protien + buffer) (in ul): & Everything in the cuvette, in \ul. \\
Enter volume of protein in cuvette (ul): & How much of your sample is in the cuvette, in \ul. \\
Enter protein's extinction coefficient: & The absorbance of a 1\,mg/ml solution in a 1\,cm cuvette
  (the ``Absorbance (at 280nm at 1 mg/ml)'' value from options 6 and 7). \\
\end{prompts}

ProtMate prints the dilution factor (total cuvette volume divided by protein volume), the
concentration in mg/ml and the total amount of protein in mg. A full example is in the Quick start
chapter.

\menuoption{4. Protein Concentration Using A280 (Multiple Readings)}
\typed{\cmd{4}, \cmd{mprot}}

The same calculation for several samples that share the same cuvette set-up and extinction
coefficient, such as fractions from a column.

\begin{prompts}
Enter total volume of liquid in cuvette (protien + buffer) (in ul): & As for option 3. \\
Enter volume of protein in cuvette (ul): & As for option 3. \\
Enter protein's extinction coefficient: & As for option 3. \\
Enter number of A280 readings: & How many samples (1 to 100{,}000). \\
Enter A280 for sample \emph{n}: & Each sample's absorbance \ldots \\
Enter yield volume for your sample \emph{n} (mL): & \ldots{} and its volume in ml. \\
\end{prompts}

ProtMate prints a table with each sample's A280, dilution factor, concentration (mg/ml) and yield
(mg), and then offers to save it (\cmd{Do you want to save your data (yes/no)?}). The saved file is
a tab-separated table (Section~\ref{sec:reports}).

\menuoption{5. Protein Concentration Using A280 (Multiple Readings, auto input)}
\typed{\cmd{5}, \cmd{mprot\_auto}}

Like option 4, but reads the readings from a text file instead of asking for each one. Prepare the
file first (Section~\ref{sec:readings-file}): each sample is an A280 value followed by its volume
in ml.

\begin{prompts}
Enter total volume of liquid in cuvette (protien + buffer) (in ul): & As for option 3. \\
Enter volume of protein in cuvette (ul): & As for option 3. \\
Enter protein's extinction coefficient: & As for option 3. \\
Enter name of input file: & The readings file, e.g. \cmd{readings.txt}. \\
\end{prompts}

\begin{session}
alice% > \inp{5}

Enter total volume of liquid in cuvette (protien + buffer) (in ul): \inp{100}
Enter volume of protein in cuvette (ul): \inp{10}
Enter protein's extinction coefficient: \inp{0.92}

Enter name of input file: \inp{readings.txt}


Sample #        A280    Dil.    [Protein] (mg/ml)       Yield (mg)
1               0.42    10      4.57                    11.425
2               0.38    10      4.13                    8.26
3               0.21    10      2.28                    6.84

Do you want to save your data (yes/no)? \inp{yes}
Enter name of output file: \inp{yields.txt}

Press the Return key to continue....
\end{session}

Up to 100{,}000 readings are used; if the file has more, ProtMate says so and uses the first
100{,}000.

\menuoption{6. Calculate Protein Extinction Coefficient}
\typed{\cmd{6}}

Calculates a protein's extinction coefficient from its numbers of tyrosine, tryptophan and cysteine
residues. (To start from a sequence instead, use option 7.)

\begin{prompts}
Enter number of Tyrosines in protein: & A whole number. \\
Enter number of Tryptophans in protein: & A whole number. \\
Enter number of Cysteines in protein: & A whole number. \\
Enter molecular weight of protein (in daltons): & The molecular weight, e.g. \cmd{16146.5}. \\
\end{prompts}

ProtMate prints the molar extinction coefficient of the native protein in water and the absorbance
at 280\,nm of a 1\,mg/ml solution in a 1\,cm cuvette. It then offers to save both values to a file.

\begin{note}
After the save question, the screen is cleared and the prompt returns without pausing. Note the
results, or save them, before you answer.
\end{note}

\menuoption{7. Calculate Protein Extinction Coefficient From a Sequence (*.seq) File}
\typed{\cmd{7}, \cmd{ext}}

Analyses a protein sequence stored in a text file (Section~\ref{sec:seq-files}): composition,
chemical formula, molecular weight, extinction coefficient and absorbance of a 1\,mg/ml solution.

\begin{prompts}
Enter name of sequence file: & The sequence file, e.g. \cmd{a.seq}. \\
Enter number of the first residue (1 to start at first residue;  0 to exit): & \cmd{1} for the
  whole sequence, a later number to start part-way through, or \cmd{0} to leave. \\
Enter number of the final residue (to end at the final residue, enter 0): & \cmd{0} for the end of
  the sequence, or the number of the last residue to include. \\
Do you want to save your data (yes/no)? & \cmd{yes} to save a report. \\
Enter name of output file: & The report file name. \\
Enter a name or unique identifier for your protein/peptide: & A single word (no spaces) that
  appears in the report's heading. \\
\end{prompts}

ProtMate prints the residues being analysed (70 per line), a warning for each unknown residue with
its position, the composition table, and the results:

\begin{session}
alice% > \inp{7}

Enter name of sequence file: \inp{a.seq}

Number of residues in sequence: 140

Enter number of the first residue (1 to start at first residue;  0 to exit): \inp{1}
Enter number of the final residue (to end at the final residue, enter 0): \inp{0}

Sequence to be analyzed:
MSFIDPYQHIIVEHQYSHKFTVVVLRATKVTKGAFGDMLDTPDPYVELFISTTPDSRKRTRHFNNDINPV
WNETFEFILDPNQENVLEITLMDANYVMDETLGTATFTVSSMKVGEKKEVPFIFNQVTEMVLEMSLEVCS


Warnings: 0

Your protein's Amino Acid composition is as follows:

Res     Num     %
======================
Tyr:    4       2.86%
Trp:    1       0.71%
Cys:    1       0.71%
   ... (one line for each amino acid) ...
Arg:    4       2.86%
Unk:    0       0%

Formula of protein/peptide: C(725)H(1116)O(221)N(180)S(8)

Molecular weight of protein/peptide (da): 16146.5

Extinction coefficient of the native protein in water: 11585

Absorbance (at 280nm & 1 mg/ml) assuming a path length of 1 cm: 0.717494

Do you want to save your data (yes/no)? \inp{yes}
Enter name of output file: \inp{a_report.txt}
Enter a name or unique identifier for your protein/peptide: \inp{ProteinA}

Your report was generated successfully.
\end{session}

If any residue is not one of the 20 standard one-letter codes, ProtMate prints
\cmd{Warning: Unknown residue at pos \emph{n}} for each one and a reminder to check the sequence
before using the results. Unknown residues are listed as \cmd{Unk} in the composition and are left
out of the formula and molecular weight. As with option 6, the screen is cleared after the save
questions, so note or save the results first.

\menuoption{8. Calculate Concentration of a Nucleic Acid from A260}
\typed{\cmd{8}, \cmd{dna}}

\begin{prompts}
What type of Nucleic Acid do you have (dna, rna or oligo)? & \cmd{dna} (double-stranded DNA),
  \cmd{rna} or \cmd{oligo}, in lower case. \\
Enter OD 260 for \ldots: & The absorbance at 260\,nm. \\
Is the dilution factor different from 1 (yes/no)? [NO] & Press \key{Return} for no, or type \cmd{yes}. \\
Enter Dilution factor: & Only asked if you answered \cmd{yes}; a whole number. \\
\end{prompts}

\begin{session}
alice% > \inp{8}

What type of Nucleic Acid do you have (dna, rna or oligo)? \inp{dna}

Enter OD 260 for dsDNA: \inp{0.65}
Is the dilution factor different from 1 (yes/no)? [NO] \inp{yes}
Enter Dilution factor: \inp{10}

Concentration = 325 ug dsDNA / ml
\end{session}

If the OD is below 0.1 or above 10, ProtMate warns that the result may not be accurate.

\menuoption{9. Calculate Melting Point of a Primer/Oligonucleotide}
\typed{\cmd{9}, \cmd{melt}}

First choose how to describe the primer:

\begin{session}
Enter input technique
[1] Manual [2] Automatic from .seq file [0] Return to main menu: [1]
\end{session}

Press \key{Return} (or type \cmd{1}) for manual entry, \cmd{2} to read the primer's sequence from a
file, or \cmd{0} to go back.

\subsubsection*{Manual entry}

\begin{prompts}
Enter length of primer/oligo (total number of nucleotides in primer/oligonucleotide): & The number
  of nucleotides. \\
Enter number of G+Cs in primer: & How many are G or C. If this equals the length, ProtMate asks
  you to confirm that the primer is all G and C. \\
Enter total salt concentration (in Molar) (Typical [salt] = 0.05 M): & The Na$^+$ concentration
  in M. Zero or a negative value is replaced by 0.05\,M. \\
Do you want to save your data (yes/no)? & \cmd{yes} to save the results, then a file name. \\
\end{prompts}

\begin{session}[fontsize=\footnotesize]
alice% > \inp{9}
Enter input technique
[1] Manual [2] Automatic from .seq file [0] Return to main menu: [1] \inp{1}

Enter length of primer/oligo (total number of nucleotides in primer/oligonucleotide): \inp{20}
Enter number of G+Cs in primer: \inp{10}
Enter total salt concentration (in Molar) (Typical [salt] = 0.05 M): \inp{0.05}


**** RESULTS ****

% G-C content: 50%

Melting point of primer = 51.78 C
Salt-adjusted melting point of primer = 46.65 C


Do you want to save your data (yes/no)? \inp{no}

Press the Return key to continue....
\end{session}

\subsubsection*{From a sequence file}

\begin{prompts}
Enter name of sequence file: & The file holding the primer's sequence (A, C, G and T only). The name
  may contain spaces. \\
Enter total salt concentration (in Molar) (Typical [salt] = 0.05 M): & As above. \\
Do you want to save your data (yes/no)? & \cmd{yes} to save, then a file name. \\
\end{prompts}

ProtMate prints the length, the number of each base, the G--C content and both melting points. If
the file contains any other character, it reports the first one and its position and stops. In
this mode the screen is cleared straight after the save question, so note or save the results
first.

\menuoption{10. Calculate your Body Mass Index (BMI)}
\typed{\cmd{10}, \cmd{bmi}}

\begin{prompts}
Enter height in feet and inches (ft [space] inches): & Both numbers on one line separated by a
  space, e.g. \cmd{5 9}. \\
Enter weight in pounds (lb): & Your weight. \\
Do you wish to enter more data (yes/no)? & \cmd{yes} to calculate another BMI. \\
\end{prompts}

ProtMate prints your height in metres, weight in kilograms, the BMI with its category, the
normal-weight range for your height, and the CDC reference ranges. Type \cmd{bmiinfo} at the
prompt to see the reference ranges on their own.

\menuoption{11. Factorial Calculator}
\typed{\cmd{11}, \cmd{fact}}

\begin{prompts}
Enter an integer: & A whole number from 0 to 20. \\
\end{prompts}

ProtMate prints, for example, \cmd{20! = 2432902008176640000}. Negative numbers and numbers above
20 are refused; a decimal number is rounded down.

\menuoption{12. Rename Oscillation Files}
\typed{\cmd{12}, \cmd{rname}}

Renames a series of diffraction images into the form Mosflm reads (see also the GUI description in
Chapter~\ref{ch:gui}). The current names must be a file string followed by \cmd{0} and a frame
number padded to three digits, such as \cmd{1.0001}; the new names are the new file string,
the three-digit frame number and \cmd{.osc}, such as \cmd{a\_001.osc}.

\begin{prompts}
Enter number of files: & How many images to rename (1 to 100{,}000), or \cmd{0} to leave. \\
Enter initial frame number: & The frame number of the first image, or \cmd{0} to leave. \\
Enter file string (e.g. for files 1.\#\#\#.*, enter "1."): & The start of the current names, e.g.
  \cmd{1.} \\
Enter desired file string (e.g to rename files a\_\#\#\#.osc, enter "a\_"): & The start of the new
  names, e.g. \cmd{a\_} \\
Working directory: & The folder holding the images, without a trailing slash: \cmd{.} for the
  current folder, a path such as \cmd{/xtal/xtal1}, \cmd{?} for help, or \cmd{0} to leave. The
  path may contain spaces. \\
Do you wish to proceed (y/n)? & \cmd{y} to rename the files listed. \\
\end{prompts}

\begin{session}
alice% > \inp{12}

Rename Osciallation Files
=========================

This module will rename sequential oscillation files to
a name format that can be used by Mosflm.

Enter number of files: \inp{3}
Enter initial frame number: \inp{1}
Enter file string (e.g. for files 1.###.*, enter "1."): \inp{1.}
Enter desired file string (e.g to rename files a_###.osc, enter "a_"): \inp{a_}
Working directory: \inp{.}
.

Starting Frame No.: 1   Ending Frame No.: 3

Frame   File
1       1.0001
2       1.0002
3       1.0003
Do you wish to proceed (y/n)? \inp{y}
Reading 1.0001...
Re-naming file 1.0001 to a_001.osc...
   ...
\end{session}

A file whose new name already exists is skipped with a message; ProtMate never overwrites a file.
The file strings may contain letters, digits and \cmd{. \_ - +}, and must not start with \cmd{-}.

\menuoption{13. Generate Configuration/Input File to Run AutoDock Vina}
\typed{\cmd{13}, \cmd{vina}}

Writes \cmd{conf.txt}, a configuration file for AutoDock Vina with a cubic search box, in the
current folder, and can then run Vina.

\begin{prompts}
Receptor file (in PDBQT format, with file extension): & e.g. \cmd{receptor.pdbqt}, or \cmd{0} to leave. \\
Ligand file (in PDBQT format with file extension): & e.g. \cmd{ligand.pdbqt}, or \cmd{0} to leave. \\
Enter the X co-ordinate for the grid center: & Centre of the search box, in \AA. \\
Enter the Y co-ordinate for the grid center: & \\
Enter the Z co-ordinate for the grid center: & \\
Enter the desired size for one side of the grid (in angstroms): & The edge length of the box. \\
Would you like to run AutoDock Vina with this configuration file (yes/no)? & \cmd{yes} to run Vina
  now (it must be installed; Chapter~\ref{ch:external}). \\
\end{prompts}

If \cmd{conf.txt} already exists, ProtMate asks before replacing it. When Vina runs, its log is
written to \cmd{vina\_log.txt} and the docked poses to \cmd{all.pdbqt}. File names may contain only
letters, digits and \cmd{. \_ - / +}. The file written looks like this:

\begin{listing}
receptor = receptor.pdbqt
ligand = ligand.pdbqt

out = all.pdbqt

center_x = 12.5
center_y = -4.2
center_z = 30.1

size_x = 24
size_y = 24
size_z = 24

exhaustiveness = 8
\end{listing}

\menuoption{14. Convert a PDBQT file to a PDB file}
\typed{\cmd{14}, \cmd{pdb}, \cmd{pdbqt}}

\begin{prompts}
Enter name of your PDBQT file leaving out its file extension (.pdbqt) (Leave empty or enter '0' to
exit module): & The file name \emph{without} \cmd{.pdbqt}. For \cmd{docked.pdbqt}, type
\cmd{docked}. The name may contain spaces. \\
\end{prompts}

ProtMate writes \cmd{docked.pdb} next to the input, asking first if that file already exists.

\menuoption{15. Generate XPLOR code to solvate a protein structure}
\typed{\cmd{15}, \cmd{xplor}, \cmd{overlay}}

Generates the X-PLOR/CNS scripts described in the GUI section (\emph{Solvate a Protein}).

\begin{session}
Which file(s) do you want to generate?
[1] First stage [2] Second stage [3] Both Stages [0] Return to main menu: [1]
\end{session}

\begin{prompts}
Enter name of protein's PDB coordinate file: & e.g. \cmd{protein.pdb} \\
Enter name of protein's MTF structure file : & e.g. \cmd{protein.mtf} \\
Enter name of solvent's PDB coordinate file: & The solvent shell (for stage 2: the file produced by stage 1). \\
Enter name of solvent's MTF structure file : & Its structure file. \\
Enter name of output file (.inp): & Where to save the script, e.g. \cmd{stage1.inp}. \\
\end{prompts}

With option \cmd{3}, ProtMate writes stage 1 and then asks for the new solvent files from stage 1
and a second output file for stage 2. File names may contain only letters, digits and
\cmd{. \_ - / +}. Run the saved scripts with X-PLOR or CNS.

% ---------------------------------------------------------------------------
\section{Other functions in the menu}

\begin{description}
  \item[\cmd{?} or \cmd{help}] Lists the command names (Table~\ref{tab:commands} has the complete
        list).
  \item[\cmd{ls}] Lists the files in the current folder (\cmd{ls -alp} on macOS and Linux, \cmd{dir}
        on Windows).
  \item[\cmd{pwd}] Shows your user name and the current folder.
  \item[\cmd{sys} or \cmd{\$}] Runs one operating-system command, typed at the \cmd{>} prompt that
        follows, for example \cmd{cp a.seq backup.seq}. The command runs with your own permissions
        through the system shell (\cmd{cmd.exe} on Windows). For safety this is disabled when
        ProtMate is run as root (macOS, Linux) or as an elevated administrator (Windows).
  \item[\cmd{abt}, \cmd{about}] Shows ProtMate's version, the copyright and licence notice, the
        disclaimer, the project's web address, and the operating system it is running on.
  \item[\cmd{0}, \cmd{quit}, \cmd{exit}, \cmd{stop} or \cmd{bye}] Leaves ProtMate.
\end{description}

\section{Commands not shown in the menu}

\begin{description}
  \item[\cmd{pmol} or \cmd{pmole}] Converts a mass/volume or volume percentage to a molar
        concentration. Choose \cmd{1} for \% (w/v), then enter the percentage and the formula
        weight (g/mol); or \cmd{2} for \% (v/v), then also the density (g/ml). The result is given in
        M and mM.
  \item[\cmd{bmiinfo}] Shows the CDC's BMI reference ranges.
  \item[\cmd{whoami}] Shows your user name.
  \item[\cmd{who}] Lists the users logged in to the computer (\cmd{query user} on Windows, where
        available).
  \item[\cmd{ps}] Lists running processes (\cmd{tasklist} on Windows).
  \item[\cmd{clear} or \cmd{clr}] Clears the screen.
  \item[\cmd{open} or \cmd{rasmol}] Asks for a PDB file name and opens it in RasMol, if RasMol is
        installed.
\end{description}

\begin{session}
alice% > \inp{pmol}

Please select the conversion you want to perform:
 [1] convert a mass concentration (%w/v) to a molar concentration
 [2] convert a percent volume (%v/v) to a molar concentration

> \inp{1}

Convert a mass concentration (%w/v) to molar concentration
==========================================================

Enter percentage solution (% (w/v)) for your compound: \inp{5}
Enter formula weight of compound (g/mole): \inp{58.44}

The molar concentration of your compound is: 0.86 M = 855.58 mM
\end{session}

\section{Command reference}

\begin{longtable}{@{}p{0.34\linewidth}p{0.6\linewidth}@{}}
\caption{Everything you can type at the \cmd{pmate} prompt.}\label{tab:commands}\\
\toprule
\textbf{Type} & \textbf{To} \\ \midrule
\endfirsthead
\toprule
\textbf{Type} & \textbf{To} \\ \midrule
\endhead
\cmd{mnu}, \cmd{menu} & Show the main menu \\
\cmd{1}, \cmd{wv-weight} & Mass needed for a \% (w/v) solution \\
\cmd{2}, \cmd{wv-vol} & Volume needed for a \% (w/v) solution \\
\cmd{3}, \cmd{prot} & Protein concentration, one A280 reading \\
\cmd{4}, \cmd{mprot} & Protein concentration, several readings typed in \\
\cmd{5}, \cmd{mprot\_auto} & Protein concentration, readings from a file \\
\cmd{6} & Extinction coefficient from residue counts \\
\cmd{7}, \cmd{ext} & Analyse a protein sequence file \\
\cmd{8}, \cmd{dna} & Nucleic acid concentration from A260 \\
\cmd{9}, \cmd{melt} & Primer melting temperature \\
\cmd{10}, \cmd{bmi} & Body mass index \\
\cmd{11}, \cmd{fact} & Factorial \\
\cmd{12}, \cmd{rname} & Rename oscillation files \\
\cmd{13}, \cmd{vina}, \cmd{VINA} & AutoDock Vina configuration \\
\cmd{14}, \cmd{pdb}, \cmd{pdbqt}, \cmd{PDB}, \cmd{PDBQT} & PDBQT to PDB \\
\cmd{15}, \cmd{xplor}, \cmd{overlay} & X-PLOR code to solvate a protein \\
\cmd{pmol}, \cmd{pmole} & Percent to molar concentration \\
\cmd{bmiinfo} & BMI reference ranges \\
\cmd{open}, \cmd{rasmol} & Open a PDB file in RasMol \\
\cmd{?}, \cmd{help} & List commands \\
\cmd{ls} & List files in the current folder \\
\cmd{pwd} & Show the current folder \\
\cmd{whoami} & Show your user name \\
\cmd{who} & Show logged-in users \\
\cmd{ps} & Show running processes \\
\cmd{sys}, \cmd{\$} & Run an operating-system command \\
\cmd{abt}, \cmd{about}, \cmd{ABT} & About ProtMate \\
\cmd{clear}, \cmd{clr} & Clear the screen \\
\cmd{0}, \cmd{quit}, \cmd{exit}, \cmd{stop}, \cmd{bye} & Quit \\
\bottomrule
\end{longtable}

% ===========================================================================
\chapter{File formats}\label{ch:files}

\section{Protein and DNA sequence files (.seq)}\label{sec:seq-files}

A sequence file is a plain text file containing the sequence in one-letter codes. Line breaks,
spaces and tabs are ignored, and upper and lower case are treated alike, so the sequence can be
split over as many lines as you like:

\begin{listing}
MSFIDPYQHIIVEHQYSHKFTVVVLRATKVTKGAFGDMLDTPDPYVELFISTTPDSRKRTRHFNNDINPV
WNETFEFILDPNQENVLEITLMDANYVMDETLGTATFTVSSMKVGEKKEVPFIFNQVTEMVLEMSLEVCS
\end{listing}

\begin{itemize}
  \item For proteins, the 20 standard amino acids are recognised (A C D E F G H I K L M N P Q R S T
        V W Y). Every other character, including numbers and punctuation, is counted as an unknown
        residue at that position.
  \item For primers (melting point), only A, C, G and T are allowed.
  \item The file must contain only the sequence. In particular, remove the header line
        (starting with \cmd{>}) from FASTA files, and any numbering, or these characters will be
        treated as residues.
  \item Sequences of up to 1{,}000{,}000 residues are accepted.
\end{itemize}

Create sequence files with any plain text editor (TextEdit in plain-text mode, Notepad, nano,
\ldots), not a word processor.

\section{A280 readings files}\label{sec:readings-file}

Used by menu option 5 and by \gui{Import\ldots} on the \gui{Multiple A280 Readings} page. The file
contains pairs of numbers: each sample's A280 followed by its volume in ml. Numbers are separated by
spaces, tabs or line breaks; one sample per line is the clearest:

\begin{listing}
0.42 2.5
0.38 2.0
0.21 3.0
\end{listing}

Commas are not separators, so convert CSV files exported from a spreadsheet by saving them as
tab-separated or space-separated text. Reading stops at the first entry that is not a number.

\section{Report files}\label{sec:reports}

Reports are plain text files with Unix (LF) line endings. They are identical whichever version of
ProtMate, and whichever operating system, wrote them. All modern text editors, including Notepad on
Windows~10 and later, display them correctly.

\begin{description}
  \item[Multiple readings] A tab-separated table, which opens directly in a spreadsheet program:
\begin{listing}
Sample #        A280    Dil.    [Protein] (mg/ml)       Yield (mg)
1       0.42    10      4.57    11.425
2       0.38    10      4.13    8.26
3       0.21    10      2.28    6.84
\end{listing}
  \item[Sequence analysis] A report headed \cmd{PROTMATE REPORTER} with the protein's name, the
        operator (your user name), the composition table, formula, molecular weight, extinction
        coefficient and absorbance.
  \item[Extinction coefficient and melting point] The values shown on the screen.
\end{description}

\section{Files for other programs}

\begin{description}
  \item[Vina configuration (\cmd{conf.txt})] Shown in the description of option 13.
  \item[X-PLOR scripts (\cmd{.inp})] Read the protein and solvent files with \cmd{structure @} and
        \cmd{coordinates @} statements, then select and write (stage 1, \cmd{overlay1.pdb}) or
        delete and write (stage 2, \cmd{overlay2.pdb}) the solvent molecules.
  \item[PDBQT to PDB] Each line of the PDBQT file is copied with columns 1--60 kept, columns 61--69
        removed, and columns 70--79 (the partial charge and the start of the AutoDock atom type)
        moved up to start in column 61, where the PDB format has the B-factor.
\end{description}

% ===========================================================================
\chapter{Safety features}

ProtMate is designed so that a typing mistake or an odd file name cannot damage your data:

\begin{itemize}
  \item \textbf{No file is overwritten without asking.} The terminal version asks
        \cmd{Overwrite it (yes/no)?} and the graphical version's save dialog asks before replacing
        a file. Renaming oscillation files never replaces an existing file.
  \item \textbf{Links are not followed when saving.} The terminal version refuses to write to a
        symbolic link (or, on Windows, another kind of reparse point), so a link cannot redirect a
        report onto another file.
  \item \textbf{File names are checked} wherever they are written into another program's input
        (Vina, X-PLOR) or used to rename files, so they cannot change the meaning of those files.
  \item \textbf{External programs are started directly}, not through a command shell, and only from
        folders named with a full path in \cmd{PATH}. The one deliberate exception is the terminal
        version's \cmd{sys} command, which exists to run your own commands and is disabled for root
        and administrator sessions.
  \item \textbf{Input limits.} Sequences are limited to 1{,}000{,}000 residues and reading lists to
        100{,}000 entries, so a wrong file cannot exhaust the computer's memory.
\end{itemize}

% ===========================================================================
\chapter{Troubleshooting}

\begin{description}[style=nextline,leftmargin=1.5em]
  \item[macOS says ProtMate ``cannot be opened'' or the developer cannot be verified.]
        See Section~\ref{sec:gatekeeper}.
  \item[macOS says ProtMate needs a newer version of macOS, or it quits immediately.]
        The ready-made app requires macOS 27 or later on Apple silicon. Build from source
        (Chapter~\ref{ch:source}) for other Macs.
  \item[\cmd{zsh: command not found: pmate} or \cmd{'pmate' is not recognized}.]
        The program is not in a folder on your search path. Run it with its full path, for example
        \cmd{\textasciitilde/bin/pmate}, or follow the installation steps to add its folder to
        \cmd{PATH}.
  \item[\cmd{zsh: permission denied: pmate} (macOS, Linux).] Make the file executable with
        \cmd{chmod +x pmate}.
  \item[Windows: ``The code execution cannot proceed because Qt6Core.dll was not found''.]
        \cmd{ProtMate.exe} was moved out of its folder or started from inside the zip file. Extract
        the whole zip file and start \cmd{ProtMate.exe} from the extracted folder.
  \item[Linux: ``could not load the Qt platform plugin xcb''.] Install the missing X11 library,
        for example \cmd{sudo apt install libxcb-cursor0} on Debian and Ubuntu, or run under
        Wayland.
  \item[CMake: ``Could not find a package configuration file provided by Qt6''.] Install the Qt
        development package (Chapter~\ref{ch:linux}), or tell CMake where Qt is with
        \cmd{-DCMAKE\_PREFIX\_PATH}. If your Qt is older than 6.5, build only the terminal version.
  \item[\cmd{Invalid filename ...}] The file was not found. Check the spelling and the current folder
        (type \cmd{pwd} and \cmd{ls} at the \cmd{pmate} prompt). File names are case-sensitive on
        Linux and usually on macOS.
  \item[Unexpected results after a typing mistake in the terminal version.] Text typed where a
        number is expected is read as 0. Start the tool again.
  \item[The screen clears before I can read the results.] Options 6 and 7, and option 9 with a
        sequence file, return to the prompt straight after the save question. Read the results, or
        answer \cmd{yes} to save them, before answering.
  \item[Many unknown residues in a sequence analysis.] The file probably contains a FASTA header,
        residue numbers or other text. Leave only the sequence in the file.
  \item[``AutoDock Vina was not found'' / ``Rasmol not available''.] See
        Chapter~\ref{ch:external}.
  \item[Vina fails with an error about the \cmd{--log} option.] ProtMate runs
        \cmd{vina --config conf.txt --log vina\_log.txt}. Releases of Vina that do not accept
        \cmd{--log} cannot be started this way; use the saved configuration file and run Vina yourself
        with \cmd{vina --config conf.txt}.
  \item[A rename shows ``Not found'' for every file.] Check the current file string and the initial
        frame number. ProtMate expects the current names to be the file string, \cmd{0}, and a
        three-digit frame number (\cmd{1.0001}).
\end{description}

\section*{Reporting problems}

If you find a bug or a wrong result, please report it on the project's issue tracker, \issues.
Include the ProtMate version (from \menupath{Help}{About ProtMate}, or \cmd{abt} in the terminal
version), your operating system, the values you entered and the result you expected. Contributions
are welcome as pull requests on GitHub.

% ===========================================================================
\appendix

\chapter{Formulas and constants}\label{app:formulas}

\section{Weight to volume ratio}
\[
  m\,[\text{g}] = \frac{\%(w/v)}{100}\times V\,[\text{ml}]
  \qquad\qquad
  V\,[\text{ml}] = \frac{m\,[\text{g}]}{\%(w/v)/100}
\]

\section{Protein concentration from A280}
\[
  D = \frac{V_\text{cuvette}}{V_\text{protein in cuvette}}
  \qquad
  c\,[\text{mg/ml}] = \frac{A_{280}\times D}{\varepsilon_{1\,\text{mg/ml}}}
  \qquad
  \text{amount}\,[\text{mg}] = c\times V_\text{sample}\,[\text{ml}]
\]
where $\varepsilon_{1\,\text{mg/ml}}$ is the absorbance of a 1\,mg/ml solution in a 1\,cm cell.
The concentration is rounded to two decimal places before the amount is calculated.

\section{Extinction coefficient}
Following Pace \emph{et al.} (1995), \emph{Protein Science} 4, 2411--2423:
\[
  \varepsilon_{280}\,[\text{M}^{-1}\,\text{cm}^{-1}] = 1490\,n_\text{Tyr} + 5500\,n_\text{Trp} + 125\,n_\text{Cys}
  \qquad
  A_{280}^{1\,\text{mg/ml}} = \frac{\varepsilon_{280}}{M_r}
\]
In Pace's method, the 125 term applies to cystines (disulfide-bonded pairs). The sequence analysis
counts every cysteine residue, so for a protein whose cysteines are reduced the value is slightly
high; with option 6 you can enter the number of disulfides instead.

\section{Chemical formula and molecular weight}
Each residue contributes its atoms as part of a polypeptide chain (the amino acid minus one water),
and one water is added for the chain's termini. The molecular weight uses average atomic masses:
C 12.011, H 1.008, N 14.007, O 15.999 and S 32.066\,Da. Unknown residues are not included.

\section{Nucleic acid concentration}
\[
  c\,[\mu\text{g/ml}] = A_{260}\times D\times F,\qquad
  F = 50 \text{ (dsDNA)},\ 40 \text{ (RNA)},\ 20 \text{ (oligonucleotide)}.
\]

\section{Primer melting temperature}
For a primer of $N$ nucleotides with $n_{GC}$ G or C residues:
\[
  T_m = \begin{cases}
     2\,(N - n_{GC}) + 4\,n_{GC} & N < 14\\[2pt]
     64.9 + 41\,\dfrac{n_{GC} - 16.4}{N} & N \ge 14
  \end{cases}\quad[\text{\degC}]
\]
\[
  T_m^\text{salt} = 81.5 + 16.6\log_{10}[\text{Na}^+] + 0.41\,(\%GC) - \frac{675}{N}\quad[\text{\degC}]
\]

\section{Percent to molar concentration}
\[
  M = \frac{\%(w/v)}{0.1\times FW}
  \qquad
  \%(w/v) = \rho\,[\text{g/ml}]\times\%(v/v)
\]

\section{Body mass index}
\[
  h\,[\text{m}] = \frac{30\,\text{ft} + 2.54\,\text{in}}{100}
  \qquad
  w\,[\text{kg}] = \frac{\text{lb}}{2.2}
  \qquad
  \text{BMI} = \frac{w}{h^2}
\]
(ProtMate uses 30\,cm per foot and 2.2\,lb per kilogram.) The normal-weight range shown is
$18.5\,h^2$ to $24.9\,h^2$, converted to pounds. Categories: below 18.5 underweight; 18.5--24.9
normal; 25.0--29.9 overweight; 30 and above obese.

\chapter{Worked example with the sample data}

Analysing the example sequence file \cmd{a.seq} (140 residues,
shown in Section~\ref{sec:seq-files}) gives:

\begin{center}\small
\begin{tabular}{@{}ll@{}}
\toprule
Formula & C$_{725}$H$_{1116}$O$_{221}$N$_{180}$S$_{8}$ \\
Molecular weight & 16146.5\,Da \\
Tyr, Trp, Cys & 4, 1, 1 \\
Extinction coefficient & 11585\,M$^{-1}$\,cm$^{-1}$ \\
Absorbance of 1\,mg/ml & 0.717 \\
\bottomrule
\end{tabular}
\end{center}

A solution of this protein with $A_{280} = 0.36$, measured after a 1-in-10 dilution
(10\,\ul{} of protein in 100\,\ul), therefore contains
$0.36\times 10 / 0.717 = 5.02$\,mg/ml.

\chapter{Licence, disclaimer and acknowledgements}\label{app:licence}

\section{Licence}

ProtMate is Copyright \textcopyright{} \copyrightyears{} Sanjeeva Wijeyesakere.

ProtMate is free software: you can redistribute it and/or modify it under the terms of the GNU
General Public License as published by the Free Software Foundation, \gplgrant. The full text of
the licence is in the \cmd{LICENSE} file of the source code, in \cmd{LICENSE.txt} in the macOS
disk image and the Windows package, and at
\url{https://www.gnu.org/licenses/gpl-3.0.html}.

In brief, the GPL allows you to run ProtMate for any purpose, to study and modify its source code,
and to distribute copies of the original or of modified versions. If you distribute ProtMate or a
modified version, in source or compiled form, you must license it under the GPL as well, keep the
copyright and licence notices, and make the corresponding source code available. The source code
of every release is available from \repo. This summary is for convenience only; the licence text
itself governs.

\section{Disclaimer}

ProtMate is a tool to assist with routine laboratory calculations. It is not a substitute for
professional judgement, and its results have not been validated for any particular use. The author
assumes no liability for the accuracy of the calculations, results or reports produced by ProtMate,
or for the functionality of the program, including any loss of samples, materials, data or time
arising from its use. You are responsible for verifying every result independently before relying
on it. The body mass index calculator is included for convenience only and is not medical advice.

The GNU General Public License states:

\begin{quote}\small
\textbf{15. Disclaimer of Warranty.} THERE IS NO WARRANTY FOR THE PROGRAM, TO THE EXTENT PERMITTED
BY APPLICABLE LAW. EXCEPT WHEN OTHERWISE STATED IN WRITING THE COPYRIGHT HOLDERS AND/OR OTHER
PARTIES PROVIDE THE PROGRAM ``AS IS'' WITHOUT WARRANTY OF ANY KIND, EITHER EXPRESSED OR IMPLIED,
INCLUDING, BUT NOT LIMITED TO, THE IMPLIED WARRANTIES OF MERCHANTABILITY AND FITNESS FOR A
PARTICULAR PURPOSE. THE ENTIRE RISK AS TO THE QUALITY AND PERFORMANCE OF THE PROGRAM IS WITH YOU.
SHOULD THE PROGRAM PROVE DEFECTIVE, YOU ASSUME THE COST OF ALL NECESSARY SERVICING, REPAIR OR
CORRECTION.

\textbf{16. Limitation of Liability.} IN NO EVENT UNLESS REQUIRED BY APPLICABLE LAW OR AGREED TO IN
WRITING WILL ANY COPYRIGHT HOLDER, OR ANY OTHER PARTY WHO MODIFIES AND/OR CONVEYS THE PROGRAM AS
PERMITTED ABOVE, BE LIABLE TO YOU FOR DAMAGES, INCLUDING ANY GENERAL, SPECIAL, INCIDENTAL OR
CONSEQUENTIAL DAMAGES ARISING OUT OF THE USE OR INABILITY TO USE THE PROGRAM (INCLUDING BUT NOT
LIMITED TO LOSS OF DATA OR DATA BEING RENDERED INACCURATE OR LOSSES SUSTAINED BY YOU OR THIRD
PARTIES OR A FAILURE OF THE PROGRAM TO OPERATE WITH ANY OTHER PROGRAMS), EVEN IF SUCH HOLDER OR
OTHER PARTY HAS BEEN ADVISED OF THE POSSIBILITY OF SUCH DAMAGES.
\end{quote}

\section{Third-party software}

The graphical version uses the Qt toolkit (\url{https://www.qt.io}), \textcopyright{} The Qt Company,
under the GNU Lesser General Public License version~3. The Windows package also includes the
\cmd{libc++} and \cmd{libunwind} runtime libraries from the LLVM project, under the Apache License
2.0 with LLVM Exceptions. The licence texts are in the \cmd{licenses} folder of the Windows
package. Qt's source code is available from \url{https://download.qt.io}. Both licences are
compatible with the GPLv3.

\section{Acknowledgements}

AutoDock Vina: O.~Trott and A.~J.~Olson (2010), AutoDock Vina: improving the speed and accuracy of
docking with a new scoring function, efficient optimization and multithreading,
\emph{Journal of Computational Chemistry} 31, 455--461.

\end{document}
